\section{Completion for every prefactor bounded below criticality}
\label{sec:all-subcritical-prefactors}
We now complete Theorem~\ref{thm:main} throughout its stated range. The preceding elementary operator-energy argument covered $Ns\le c_0$. For the remaining fixed prefactors, the appropriate depth cost is the covariance action $\kappa(\nu)H$, rather than the mean-shift energy $H$.

Fix $\nu_{\max}<1$. If $\nu_{\max}\le c_0$, Section~\ref{sec:assembly} already proves the theorem. Otherwise restrict temporarily to $\nu=Ns\in[c_0,\nu_{\max}]$. This is a compact subset of $(0,1)$, and $\kappa_{\min}:=\inf\kappa(\nu)>0$ on it. Fix a small $t>0$ and define
\[
 q=Q_{H,N,\nu}(t,n),\qquad K=(1+\delta)\kappa(\nu)H.
\]
For $n=4\cdot2^{\lfloor\eta H\rfloor}$ and $N\ge e^{CH}$,
\[
 q/H=\kappa(\nu)+o(1)
\]
uniformly on this compact interval. Choose $\delta$ sufficiently small and then the depth threshold sufficiently large that
\[
 0<q\le K,\qquad 1-q/K\le2^{-18}.
\]
Theorem~\ref{thm:subcritical-prefactor-memorizer} gives a comparison fitting the bounded empirical atom means with expected loss at most $r_0+2^{-20}$ and normalized entropy at most $K$. Its constants and width threshold are uniform on the compact interval, after shrinking $\eta$ and increasing $C$ if needed.

For the actual posterior let $A=\{\|z_w(S)\|_n>t\}$ and use the nonnegative auxiliary energy $Q(w)=q\mathbf1_A$. The event-entropy floor and the zero and memorizer comparisons imply
\[
 \E_\Pi[\widehat R_S(z_w)+\lambda Q(w)]
 \le\min\{r_0+D_S,r_0+2^{-20}+\lambda K\}
                       +s+2\lambda s\log2.
\]
The auxiliary energy has exactly the floor required by Lemma~\ref{lem:segment}. For every fixed $\Lambda$, the error $s+2\lambda s\log2$ tends uniformly to zero over $\lambda\le\Lambda/H$ and the prescribed widths. Take $t\le2^{-24}$ and increase the depth and width thresholds so that this error is at most $2^{-20}$. The segment distance is then less than $1/256$, exactly as in Section~\ref{sec:assembly}. Because $|Y_i|<1$, clipping each draw can only reduce its loss; the same auxiliary energy still has the required floor for the clipped norm. Thus all stated raw and clipped means satisfy the same segment estimate.

For $\lambda\ge\Lambda/H$, Lemma~\ref{lem:subcritical-high-temperature} gives
\[
 \|\E_\Pi z_w(S)\|_n\le t+
 \sqrt{5\left\{\frac{1+s}{\lambda q}+\frac{2s\log2}{q}\right\}}.
\]
The identical upper bound holds for the mean of clipped draws by splitting on $A$ and using that clipping decreases the empirical norm. Choose $t$ smaller if necessary, then choose a fixed $\Lambda$ sufficiently large in terms of $\kappa_{\min}$, and finally increase the depth threshold. The displayed norm is at most $2^{-15}$ throughout the high-temperature range. Clipping a mean also decreases that norm.

The low- and high-temperature branches therefore have exactly the common witness and segment properties used in Section~\ref{sec:assembly}. Its population conversion gives excess at least $2^{-16}$, including mixtures and clipping after mixing. The sample event depends only on the observations, so it is common to the whole compact prefactor interval and to the already established $\nu\le c_0$ range. Taking the smaller support-growth constant and larger thresholds from the two constructions proves Theorem~\ref{thm:main} for all $0<Ns\le\nu_{\max}$.

The same replacement extends the Gaussian-noise conclusion of Theorem~\ref{thm:gaussian-fixed-accuracy}. That proof already has $\widehat R(0),D_S,\|b-2g\|_n^2\le1$ on its sample event. Its comparison clips empirical atom means to the fixed interval $[-4,4]$, for which Theorem~\ref{thm:subcritical-prefactor-memorizer} provides the same prescribed RMS error and cost ratio. No other step changes. Consequently its raw-mean risk lower bound $1/16+2^{-21}$ holds uniformly over every fixed range $0<Ns\le\nu_{\max}<1$, with constants allowed to depend on $\nu_{\max}$.

Finally the bounded-output-likelihood conclusion of Corollary~\ref{cor:bounded-output-likelihood} also extends to this full range. For either output map $\psi$, the event $\{\|\psi(z(S))\|_n>t\}$ is contained in the prior event used by the entropy floor. Oddness preserves the centered zero comparison, and contraction preserves its second-moment bound. The comparison fits the same bounded empirical means for clipping and their bounded inverse-tanh transforms for a final tanh. Thus the segment argument and the high-temperature event estimate apply directly to the actual prediction $\psi(z)$. This extension concerns the corresponding bounded-output likelihood, not a claim that raw affine conjugacy survives such an output map.
